% 日本語論文のため LuaTeX + luatexja を用いる。pdflatex では日本語文字ごとに
% "LaTeX Error: Unicode character" が発生し、本文が PDF から欠落する。
% 同梱の英語会議用テンプレート (mdpi / iclr2024 / agents4science_2025) は
% CJK 非対応のため使用せず、自前のプリアンブルで組む。
\documentclass[11pt,a4paper]{article}
\usepackage{luatexja-fontspec}
\setmainjfont{IPAexMincho}
\setsansjfont{IPAexGothic}
\usepackage{graphicx}
\usepackage{booktabs}
\usepackage{amsmath}
% \mathbb は amssymb 由来。amsmath だけでは未定義になり、数式が組めない。
\usepackage{amssymb}
\usepackage[margin=25mm]{geometry}
\usepackage[hidelinks]{hyperref}
\usepackage{caption}
\captionsetup{font=small}

\title{推論専用公開モデルに対する分散ファインチューニングの実行正当性検証\\
\large ---- OpenDDE 拡散モジュールを 4 ノード 16 GPU で 30 エポック学習した記録 ----}
\author{松澤 拓海}
\date{2026年8月15日}

\begin{document}
\maketitle

\begin{abstract}
マルチノード分散学習では、計算資源が確保されているにもかかわらず一部のランクしか
計算に参加していないという故障が起こりうる。この故障は例外を送出せず、ジョブスケジューラ上でも
正常なジョブと区別がつかないため、通常の実行ログからは検出できない。
本研究は、全ランクが自身のホスト名と GPU の一意識別子を報告し、それを NCCL の集団通信で
集約するという検証機構を実装し、その集約結果の出力自体をノード間通信成立の証拠として用いる。
検証対象として、推論専用ビルドとして公開されている生体分子構造予測モデル OpenDDE を選び、
技術報告書に記載されたハイパーパラメータから学習ループを再構成した。
理化学研究所 RIKYU クラスタの 4 ノード 16 GPU (NVIDIA GB200) 上で 30 エポック、
ランクあたり 60 回の勾配更新 (全ランク合計 960 ステップ) を 1411 秒 (23.5 分) で実行し、
16 ランクすべてが 4 ノードに均等配置され、同数のステップを実行したことを確認した。
さらに、1 ノードで 4 ノードを要求する負の対照実験を行い、検証機構が実際に
FAIL 判定と非ゼロ終了を返すことを確認した。
一方、FoldBench の抗体--抗原インターフェース予測における DockQ 成功率は
事前学習済みモデルと微調整済みモデルで同一 (ともに 10.53\%、19 インターフェース中 2 件) であり、
平均 DockQ は 0.0844 から 0.0693 へ低下した (対応のある差の平均 $-0.0151$、
Wilcoxon 符号順位検定 $p = 0.16$)。学習損失の回帰直線の傾きは
$-6.8\times10^{-5}$/エポック ($r^2 = 0.018$、$p = 0.49$) であり、統計的に有意な減少は検出されなかった。
これらは実験前に文書化した予測と一致する。本研究の主張は精度向上ではなく、
分散実行の正当性を検証可能にする手法の提示にある。
\end{abstract}

\section{はじめに}

大規模モデルの学習では複数ノードにまたがる分散実行が前提となる。
その際に問題となるのは、計算が明示的に失敗する場合ではなく、
\textbf{失敗しないまま誤った結果を出す}場合である。

具体例を挙げる。分散学習の実行基盤が環境変数 \texttt{RANK} や \texttt{WORLD\_SIZE} を
静的な単一値として全プロセスに配布した場合、16 個のプロセスはいずれも
「自分は単独プロセスである」と認識する。勾配同期は行われず、実質的に 16 個の独立した学習が
走るが、例外は送出されず、ジョブは正常終了し、スケジューラ上は 4 ノードを確保した
正常なジョブとして記録される。実行ログを開いても、学習ステップと損失が出力されている以上、
異常を示す手がかりは存在しない。

本研究はこの種の故障を検出可能にすることを目的とする。
中心となる考え方は単純である。全ランクに自身のホスト名と GPU の一意識別子を報告させ、
その報告を集団通信で集約する。\textbf{集約結果が出力されたという事実自体が、
ノード間でデータが実際に流れたことの証拠となる}。
確保されたノード数ではなく、報告に現れた相異なるホスト名の個数を数えることで、
「確保されている」と「計算している」を分離できる。

検証対象として OpenDDE~\cite{project-2026-folding} を用いた。
これは Apache-2.0 で公開された全原子生体分子基盤モデルであり、
抗体--抗原共フォールディングにおいて既存のオープンモデルを上回る精度が報告されている。
選定理由は後述するが、公開版が推論専用ビルドであるため、
分散学習を行うには学習ループ自体を外部から構成する必要があり、
検証機構を組み込む余地があった点が大きい。

\section{関連研究}

\subsection{全原子構造予測モデル}

AlphaFold3~\cite{abramson-2024-accurate} は、タンパク質・核酸・低分子・イオンを含む
複合体の原子レベル同時モデリングを実現した。その拡散モジュールは、
ノイズを加えた座標から真の座標を復元する学習によって訓練される。
OpenDDE~\cite{project-2026-folding} はこの枠組みを踏襲しつつ Pairformer の隠れ次元を
128 から 384 に拡大し、原子レベルの潜在推論ブロックを導入したモデルである。

\subsection{評価ベンチマーク}

FoldBench~\cite{xu-2025-benchmarking} は 1,522 個の生物学的アセンブリを
9 種のタスクに分類したベンチマークである。本研究では抗体--抗原インターフェース予測を用いる。
インターフェース予測の品質は DockQ~\cite{basu-2016-dockq} で測る。
DockQ は CAPRI の分類基準を連続値化した指標で、0.23 を超えると acceptable、
0.49 を超えると medium、0.80 を超えると high と分類される。
算出には OpenStructure~\cite{studer-2020-modeling} の
\texttt{compare-structures} を用いた。

\subsection{分散学習の検証}

PyTorch の DistributedDataParallel~\cite{li-2020-pytorch} は、
各プロセスがモデルの複製を保持し、逆伝播のたびに勾配を集団通信で平均する方式をとる。
本研究が扱う故障は、この枠組みが正しく初期化されなかった場合に生じる。
分散学習の性能評価は広く行われている一方、
\textbf{実行が意図した並列度で行われたことの事後検証}に焦点を当てた報告は少ない。

\section{背景}

\subsection{公開版 OpenDDE には学習コードが存在しない}

OpenDDE の技術報告書は「学習コード、推論パイプライン、チェックポイント、
ベンチマークを公開する」と述べている~\cite{project-2026-folding}。
しかし本研究の時点で公開リポジトリを調査した結果、
学習ループ、損失関数、オプティマイザのいずれも存在しなかった。
ドキュメント内にも「推論専用ビルド (inference-only build)」という記述があり、
公開されているのは推論経路のみである。

したがって本研究における学習は、公開コードの再現ではなく、
技術報告書に記載されたハイパーパラメータからの再構成である。この区別は重要であり、
本稿で報告する学習の挙動は OpenDDE 開発陣の学習手順と一致する保証がない。

\subsection{再構成に利用できた要素}

一方で、推論専用ビルドにも学習に転用できる要素が残されていた。
第一に、拡散モジュールの \texttt{forward} が
「ノイズ付き座標とノイズ強度を受け取り、デノイズ済み座標を返す」
単一ステップのデノイザとして公開されている。これは拡散学習の損失を接続する対象そのものである。
第二に、実装内の in-place 最適化はすべて勾配が有効な場合に自動的に無効化される設計になっており、
逆伝播を妨げない。第三に、カスタム CUDA LayerNorm は \texttt{backward} を実装した
\texttt{torch.autograd.Function} として定義されている。

ただしトランク (入力埋め込み、MSA モジュール、テンプレート埋め込み、Pairformer) は
\texttt{torch.set\_grad\_enabled(False)} で囲まれており、勾配を受け取れない。
本研究ではこれを制約ではなく設計として受け入れ、拡散モジュールのみを学習対象とした。

\section{手法}

\subsection{学習目的関数の再構成}

拡散学習の損失は EDM~\cite{karras-2022-elucidating} の定式化に従う。
ノイズ強度 $\sigma$ を対数正規分布から標本化する:
\begin{equation}
\sigma = \sigma_{\text{data}} \exp(p_{\text{mean}} + p_{\text{std}} \cdot n),
\quad n \sim \mathcal{N}(0, 1)
\end{equation}
ここで $p_{\text{mean}} = -1.2$、$p_{\text{std}} = 1.5$、$\sigma_{\text{data}} = 16.0$ は
OpenDDE 技術報告書 Table 1(b) に記載された値である~\cite{project-2026-folding}。

正解座標 $x_{\text{GT}}$ に対しノイズを加え、デノイザ $D_\theta$ で復元し、
AlphaFold3 補遺 3.7.1 の重み付き平均二乗誤差~\cite{abramson-2024-accurate} を損失とする:
\begin{equation}
\mathcal{L} = \mathbb{E}_\sigma \left[
\frac{\sigma^2 + \sigma_{\text{data}}^2}{(\sigma \cdot \sigma_{\text{data}})^2}
\cdot \frac{1}{3N}\sum_{i} w_i \left\| D_\theta(x_{\text{noisy}}, \sigma)_i
- \hat{x}_{\text{GT},i} \right\|^2 \right]
\end{equation}
$\hat{x}_{\text{GT}}$ は予測構造に対して重み付き剛体重ね合わせを施した正解座標であり、
剛体運動の自由度を損失から除去する。$w_i$ は原子種別ごとの重みである。

\subsection{正解座標のモデル原子順序への整列}

推論専用ビルドは正解座標を生成する経路を持たない。特徴量化が出力するのは
CCD の理想化配座であって実験座標ではなく、ラベル用のスキーマは定義だけが残された
未使用の状態にある。したがって教師信号の構成は自前で行う必要がある。

本研究では、各エントリに同梱される FASTA (鎖ごとの配列) と NPZ (実験座標) から
教師信号を構成した。FASTA のヘッダに記された鎖 ID をそのまま入力 JSON に指定することで、
モデルが構築する原子配列と正解座標側の鎖命名が構成上一致する。
両者を \texttt{(鎖 ID, 残基番号, 原子名)} の三つ組で結合し、
一致しない原子および NaN として格納されている未観測原子はマスクした。

なお、同梱される mmCIF は使用していない。前処理済み mmCIF には
\texttt{\_entity\_poly} カテゴリが欠落しており、OpenDDE の mmCIF パーサが
全ポリマーをリガンドと誤認して原子単位に分解するためである。
この経路では結合率が 0\% となった。

整列の正しさは幾何的に検証した。8 エントリにおいて、隣接原子間距離の中央値は
すべて 1.52--1.53 \AA{} であった。これは C--C および C--N 共有結合長に一致し、
原子が正しい順序で対応付いていなければ得られない値である。

\subsection{分散実行の検証機構}

各ランクは学習開始前に、自身のホスト名、GPU の製品名、GPU の一意識別子 (UUID) を
1 行の JSON として出力する。続いてこれらのレコードを \texttt{all\_gather\_object} で
集約し、ランク 0 が集計結果を出力する。集計には次の 3 つが含まれる:
相異なるホスト名の個数、相異なる GPU UUID の個数、ホストごとのランク数である。

期待ノード数を設定として与え、集計されたホスト数がこれを下回る場合は
\texttt{FULL\_VALIDATION: FAIL} を出力し、プロセスの終了コードを 1 とする。
終了コードまで非ゼロにするのは、FAIL と出力しながら 0 で終了した場合、
スケジューラ側には成功として記録され、この検証機構が取り除こうとしている混同が
そのまま残るためである。本稿の実装では当初この点が誤っており、
FAIL を出力しながら終了コード 0 を返していた。この誤りは自己査読の過程で発見し、
後述の負の対照実験 (第~\ref{sec:negcontrol} 節) の実施前に修正した。
したがって第~\ref{sec:negcontrol} 節の対照実験は修正後のコードで、
第~\ref{sec:placement} 節の 4 ノード実行は修正前のコードで行われている。
判定が PASS の場合の終了コードは修正の前後で 0 のまま変わらないため、
報告した 4 ノード実行の結果はこの修正の影響を受けない。

データ分割については、各ランクは自身の担当断片のみを計算する。
ただし\textbf{断片は互いに素ではない}。ランク間でステップ数を揃えるため、
標本数がランク数の倍数でない場合は先頭の標本を巡回して埋めており、
本実験では 27 標本を 16 ランクに配る際に 5 標本が 1 エポック内で 2 回現れている。
ステップ数を揃える必要があるのは、ランクごとにステップ数が異なると
集団通信の回数が食い違い、実行が停止するためである
(実際、この不整合は実装過程で NCCL のウォッチドッグによる停止として顕在化した)。
したがってここで排除できているのは「一部のランクだけが計算している」状態であって、
「全ランクが完全に異なるデータを見ている」ことではない。

\section{実験設定}

\subsection{計算環境}

理化学研究所 RIKYU クラスタを用いた。1 ノードあたり NVIDIA GB200 を 4 基、
CPU 144 コア、メモリ 1.72 TB を備え、アーキテクチャは aarch64、
ジョブスケジューラは Slurm である。実行は 4 ノード 16 GPU を占有して行った。

\subsection{学習データ}

OpenFold3 の PDB 学習サブセット (1,000 エントリ) から、
タンパク質のみで構成され、全鎖が異なる配列を持ち、非標準残基を含まない
2 鎖複合体を選定した。全鎖異配列を条件としたのは、
同一配列の鎖が複数ある場合に鎖の対応付けが一意に定まらず、
OpenDDE が対称性を扱う特徴量を持たない以上、
任意のラベル付けに向けてモデルを学習させることになるためである。

この条件を満たすエントリはサブセット全体で 54 件存在した。
そのうち\textbf{解像度の昇順で先頭 30 件}を前処理の対象とした。
前処理では実験座標が解決している原子の割合が 0.5 未満のエントリを除外する規則を設けたが、
この 30 件では該当が 0 件であった。
学習時にさらに、モデルが構築する原子数が 4,000 を超える 3 件 (最大 6,150 原子) を除外し、
最終的に 27 件を用いた。原子数の上限を設けない場合、
構造トークンのリファイナが単一のテンソルに 144 GiB を要求し、
実行が確保したランクの担当データに依存して成否が変わる。

選定を解像度順としたことには副作用がある。得られた 27 件は
サブセット中で最も実験的に良く決まった構造であり、無作為標本ではない。
学習データの分布としては偏っている。

なお、これらのエントリの登録日は 1976 年から 2021 年 9 月であり、
OpenDDE の学習データカットオフ (2021 年 9 月) の内側にある。
すなわち\textbf{モデルにとって既知の分布であり、新しい情報を含まない}。

\subsection{学習設定}

学習率 $1\times10^{-4}$ の AdamW、勾配クリッピング閾値 10.0、
トランクの再循環回数 4、1 ステップあたりのノイズ標本数 4 とし、30 エポック実行した。
学習対象は拡散モジュールの 203.3 M パラメータであり、
モデル全体 655.8 M のうち残りは凍結した。

\subsection{評価設定}

FoldBench の抗体--抗原インターフェースから 16 ターゲットを対象とし、
原子数 8,000 を超える 2 件を除外して 14 ターゲット (19 インターフェース) を採点した。
なお原子数の定義は学習時と推論時で異なる。学習時の 4,000 はモデルが構築する
原子配列の長さ、推論時の 8,000 は正解 mmCIF の \texttt{ATOM}/\texttt{HETATM} 行数であり、
両者は直接比較できない。いずれも計算資源上の上限として個別に決めたものである。

拡散サンプリングは 20 ステップ、1 ターゲットあたり 1 サンプル、乱数種は 42 とした。
20 ステップという値は既定値をそのまま用いたもので、
サンプリング品質が飽和する点を測定して選んだわけではない。
また\textbf{乱数種は 1 つのみであり、拡散サンプリング由来の分散を推定していない}。
本稿で報告する 2 モデル間の差は、この分散と比較されていない点に注意が必要である。

MSA およびテンプレート検索は\textbf{学習時・推論時のいずれにおいても行っていない}。
学習時はダミーの MSA・テンプレート特徴量を与えており、
検索を伴う本来の入力条件とは異なる。推論時に検索を省いていることは
絶対精度に対する大きな不利であるが、事前学習済みモデルと微調整済みモデルの双方に
同一条件を課すため、両者の比較においては相殺される。
したがって\textbf{本稿の絶対値を FoldBench の公開リーダーボードと比較してはならない}。

\section{結果}

\subsection{分散実行は正しく行われた}
\label{sec:placement}

30 エポックの学習は 1411 秒 (23.5 分) で完了し、検証は成功と判定された。
集計結果を表~\ref{tab:placement} に示す。
なおジョブ全体の実時間は 2459 秒 (41.0 分) であるが、
その差はコンテナイメージの構築とデータのキャッシュ読み込みに費やされており、
学習ループの実行時間ではない。

\begin{table}[htbp]
\centering
\caption{分散実行の集計結果。集計自体が \texttt{all\_gather\_object} の出力であり、
その存在がノード間通信の成立を示す。}
\label{tab:placement}
\begin{tabular}{ll}
\toprule
項目 & 値 \\
\midrule
相異なるホスト数 & 4 (c050, c054, c055, c056) \\
ホストごとのランク数 & 各 4 \\
相異なる GPU UUID 数 & 16 \\
総ランク数 & 16 \\
ランクあたりの勾配更新回数 & 60 (= 2 ステップ $\times$ 30 エポック) \\
全ランク合計のステップ数 & 960 (= 16 ランク $\times$ 60) \\
学習ループの実行時間 & 1411 秒 (23.5 分) \\
判定 & PASS \\
\bottomrule
\end{tabular}
\end{table}

GPU UUID が 16 個すべて相異なることは、物理的に別々の 16 基の GPU が
使用されたことを意味する。ランクあたりの勾配更新回数が 60 で均等であることは、
全ランクが同量の計算を担ったことを示す。
ここで\textbf{960 はランク数とステップ数の積であり、
モデルが受けた勾配更新の回数は 60 である}。DDP は全ランクの勾配を平均して
1 回の更新を行うため、960 を更新回数と読むことはできない。
この区別は、27 標本のデータセットに対して行われた学習が
どの程度の規模であったかを判断する上で重要である。

学習中、勾配ノルムは全ランクで一致していた (エポック 3 で 0.056、エポック 14 で 0.028)。
これは DDP が勾配を平均していることの直接的な帰結であり、
損失がランクごとに異なる (各ランクが異なる構造を担当するため) のと対照的である。

\subsection{学習損失に有意な減少は検出されなかった}

図~\ref{fig:loss} にエポックごとの損失推移を示す。
値は 0.0373 から 0.0556 の範囲を上下するのみで、
初回エポックの 0.0498 に対し最終エポックは 0.0469 である。
エポック番号に対する線形回帰の傾きは $-6.78\times10^{-5}$/エポックで、
$r^2 = 0.018$、$t = -0.71$ ($p = 0.49$)、傾きの 95\% 信頼区間は
$[-2.6\times10^{-4},\; +1.3\times10^{-4}]$ である。
すなわち\textbf{減少は有意ではない}。

ここで主張できるのは「傾向が検出されなかった」ことであって、
「傾向が存在しない」ことではない。信頼区間の下端は 30 エポックで 0.0077 の減少に相当し、
これは観測された変動幅 (0.018) の半分近くを占める。
本実験の統計的検出力では、この程度の減少があったとしても区別がつかない。

\begin{figure}[htbp]
\centering
\includegraphics[width=0.85\textwidth]{images/chart/training_loss.pdf}
\caption{30 エポックの学習損失推移。全ランク平均。破線は線形回帰
(傾き $-6.78\times10^{-5}$/エポック、$r^2 = 0.018$、$p = 0.49$)。}
\label{fig:loss}
\end{figure}

\subsection{DockQ は改善しなかった}

FoldBench 抗体--抗原インターフェースにおける評価結果を
表~\ref{tab:dockq} に示す。

\begin{table}[htbp]
\centering
\caption{FoldBench 抗体--抗原インターフェース (19 インターフェース) の評価結果。
MSA・テンプレート検索なし。}
\label{tab:dockq}
\begin{tabular}{lrr}
\toprule
指標 & 事前学習済み & 微調整済み (30 エポック) \\
\midrule
DockQ 成功率 ($>0.23$) & 10.53\% (2/19) & 10.53\% (2/19) \\
DockQ 平均 & 0.0844 & 0.0693 \\
DockQ 中央値 & 0.039 & 0.034 \\
DockQ 最大値 & 0.317 & 0.300 \\
medium 率 ($>0.49$) & 0\% & 0\% \\
lDDT 平均$^{\dagger}$ & 0.660 & 0.643 \\
\bottomrule
\end{tabular}

\vspace{2mm}
\begin{minipage}{0.9\textwidth}
\footnotesize $\dagger$ lDDT は \texttt{compare-structures} が構造全体に対して返す値であり、
インターフェース単位ではない。同一ターゲットに複数のインターフェースがある場合、
同じ値が重複して集計に入っている。したがって DockQ とは独立な 19 点の標本ではない。
\end{minipage}
\end{table}

主要指標である DockQ 成功率は両者で完全に一致した。
平均 DockQ および lDDT 平均はいずれもわずかに低下している。
図~\ref{fig:cumulative} の累積分布では 2 本の曲線がほぼ重なる。
両モデルとも acceptable 閾値 (0.23) を超えるインターフェースが 2 件ずつ存在し、
最大値は事前学習済みで 0.317、微調整済みで 0.300 である。
すなわち分布は閾値の手前で終わっているのではなく、
閾値をまたいだ位置に少数の例が残り、medium 閾値 (0.49) には届いていない。

\begin{figure}[htbp]
\centering
\includegraphics[width=0.85\textwidth]{images/chart/dockq_cumulative.pdf}
\caption{DockQ の累積分布。横軸の各閾値を超えるインターフェースの割合。}
\label{fig:cumulative}
\end{figure}

インターフェースごとに見ると (図~\ref{fig:paired})、変化は両方向に散らばるが、
\textbf{件数は対称ではない}。19 インターフェースのうち改善は 6 件、悪化は 13 件であり、
対応のある差の平均は $-0.0151$ (標準偏差 0.0903) である。
Wilcoxon 符号順位検定では $p = 0.16$、対応のある $t$ 検定では
$t(18) = -0.73$ ($p = 0.48$) であり、いずれも有意水準 5\% では有意でない。
改善幅が大きい例 (8fdo: 0.068 $\to$ 0.300、8hxq: 0.176 $\to$ 0.283) と
悪化幅が大きい例 (8fdd-1: 0.242 $\to$ 0.058、8ux6-1: 0.317 $\to$ 0.174) が
併存するため平均差が小さくなっているのであって、
改善と悪化が同数だからではない。

成功率が両者で 10.53\% と一致している点についても、内訳は一致していない。
閾値を超えたインターフェースは事前学習済みが \{8fdd-1, 8ux6-1\}、
微調整済みが \{8fdo, 8hxq\} であり、\textbf{両者に重なりはない}。
成功率という要約値の一致は、同じターゲットで同じ性能が出たことを意味しない。

\begin{figure}[htbp]
\centering
\includegraphics[width=0.8\textwidth]{images/chart/dockq_per_interface.pdf}
\caption{インターフェースごとの DockQ。各行が 1 インターフェースで、
両モデルの値を線で結んでいる。}
\label{fig:paired}
\end{figure}

\subsection{負の対照実験}
\label{sec:negcontrol}

前節までの結果は、検証機構が PASS を返した実行のみを示している。
PASS を返す機構は、常に PASS を返す機構と区別がつかない。
そこで、同一のコミットに対し\textbf{1 ノード 4 GPU のみを確保した上で
期待ノード数 4 を要求する}実行を行った。

結果を表~\ref{tab:negcontrol} に示す。検証機構は
\texttt{PILOT\_VALIDATION: FAIL reason=expected\_4\_nodes\_saw\_1} を出力し、
ランク 0 は終了コード 1 で終了した。これを受けて \texttt{srun} が
ジョブステップ全体を打ち切り、実行基盤側にも失敗として記録された。

\begin{table}[htbp]
\centering
\caption{負の対照実験。1 ノード 4 GPU で 4 ノードを要求した場合の集計結果。
学習自体は正常に完走している点に注意されたい。}
\label{tab:negcontrol}
\begin{tabular}{lll}
\toprule
項目 & 本実験 (4 ノード) & 負の対照 (1 ノード) \\
\midrule
確保した GPU 数 & 16 & 4 \\
相異なるホスト数 & 4 & 1 (c183) \\
相異なる GPU UUID 数 & 16 & 4 \\
期待ノード数 & 4 & 4 \\
判定 & PASS & \textbf{FAIL} \\
終了コード & 0 & \textbf{1} \\
実行基盤の記録 & completed & \textbf{failed} \\
\bottomrule
\end{tabular}
\end{table}

この対照実験で重要なのは、\textbf{学習そのものは何の問題もなく完走している}ことである。
8 ステップが実行され、損失は 0.0503 から 0.0781 へ推移し、
損失曲線・配置記録・チェックポイントがいずれも正常に書き出された。
実行時間は 47.4 秒であった。
すなわち「エラーが出ていないこと」「成果物が生成されていること」「損失が記録されていること」は、
いずれも意図した並列度で実行されたことの証拠にならない。
この実行を失敗と判定した根拠は、集約されたホスト数が 1 であったという一点のみである。

この対照実験がなければ、本稿の検証機構に関する主張は
コードの読み下しにとどまり、実験結果ではなかった。
「1 ノードに集中していれば成功として記録されない」という記述は、
対照実験を経て初めて観測された事実になる。

\section{考察}

\subsection{予測との一致}

本研究では実験実施前に、分散実行は正しく動作する一方 DockQ は改善しないか
わずかに低下すると予測し、その理由を文書として記録した。結果は予測と一致した。
理由として挙げた「学習データが事前学習分布の内側にある」という点は、
学習損失に有意な減少が検出されなかったことと整合する。

ただしこの整合は弱い証拠である。損失が動かない理由としては、
分布内データであること以外に、27 標本という規模、ランクあたり 60 回という更新回数、
拡散モジュールのみを学習対象としたこと、ダミーの MSA 特徴量を与えたことなど
複数の説明が同時に成り立つ。本実験はこれらを分離していない。

予測を事前に記録したのは、結果を得てから解釈を都合よく変更することを防ぐためである。
本研究のように否定的な結果が予想される場合、この手続きの価値は特に大きい。

\subsection{エポック数を伸ばして得られたもの}

先行して実施した 3 エポックの実行 (全ランク合計 96 ステップ、166 秒) と比較すると、
本実行は 10 倍の規模である。この延長によって得られた情報は 2 つある。

第一に、\textbf{損失の推移を統計量として扱えるようになった}。
3 エポックでは回帰の当てはめ自体が意味を持たない。30 点あることで傾きと
その信頼区間を報告でき、「減少が検出されなかった」ことを検出力とともに述べられる。
ただし前述のとおり、これは傾向の不在の証明ではない。

第二に、\textbf{より長い実行時間にわたる分散実行の安定性が確認された}。
166 秒では、通信が徐々に劣化する種の異常は発現する時間すら与えられていない。
23.5 分にわたり勾配同期が維持されたことは、より強い主張である。
とはいえ 23.5 分もまた、実運用の学習ジョブと比べれば短い。
数日規模の実行で現れる異常についてはなお何も述べていない。

\section{限界}

本研究には次の限界がある。

\textbf{学習コードは公開版の再現ではない。} OpenDDE の公開リポジトリに
学習コードが存在しないため、本研究の学習ループは技術報告書の記述からの再構成である。
開発陣の実際の学習手順と一致する保証はない。

\textbf{InfiniBand を使用していない。} ジョブコンテナに \texttt{libibverbs} が
含まれないため、NCCL は IPoIB 上の TCP ソケットにフォールバックしている。
通信は成立するが、GB200 の相互接続性能は発揮されていない。
本研究は分散実行の\textbf{正当性}を扱うものであり、性能については何ら主張しない。

\textbf{スケーリングを測定していない。} 1、4、8、16 GPU での処理時間比較を
行っていないため、ノード数を増やした際の効率は不明である。

\textbf{検出できる故障の範囲が限られている。} 本機構が検出するのは
「参加しているホスト数が期待を下回る」という故障であり、
負の対照実験で確認したのもこの 1 種類のみである。
たとえば全ランクが別々のノードにありながら勾配が誤って同期される、
一部のランクだけが古いパラメータを保持する、といった故障は
ホスト数の集計では検出できない。本稿の勾配ノルム一致の観察は
これらに対する補助的な証拠にとどまり、対照実験による裏付けを持たない。

\textbf{評価規模が小さく、乱数種は 1 つである。} 19 インターフェースでの比較であり、
DockQ 成功率における 1 件の差は約 5 ポイントに相当する。
加えて拡散サンプリングを 1 種のみで実行しているため、
サンプリング由来の分散を推定していない。観測された平均差 $-0.0151$ が
その分散より大きいかどうかは判定できていない。
本研究の結論は「改善が観測されなかった」ことであり、
「改善しないことが示された」ではない。

\textbf{学習データが小規模である。} 27 件、ランクあたり 60 回の更新という規模は、
学習対象である 203.3 M パラメータに対して意味のある更新を与えるには小さい。
条件を緩めずに手元のサブセットから得られる上限は 54 件であり、
本実験ではそのうち解像度上位 30 件を起点としている。
また選定が解像度順であるため、学習データは無作為標本ではない。

\section{結論}

推論専用ビルドとして公開されたモデルに対し、技術報告書の記述から学習ループを再構成し、
4 ノード 16 GPU の分散環境で 30 エポック、ランクあたり 60 回の勾配更新を実行した。
16 ランクすべてが 4 ノードに均等配置され、相異なる 16 基の GPU 上で
同数のステップを実行したことを、集団通信による集約結果として確認した。
さらに 1 ノードで 4 ノードを要求する負の対照実験により、
この検証機構が実際に失敗を返すことを確認した。

精度は改善しなかった。DockQ 成功率は事前学習済みと同一であり、
平均 DockQ は 0.0151 低下した (Wilcoxon 検定 $p = 0.16$、有意でない)。
学習損失にも有意な減少は検出されなかった。
これらは実験前に記録した予測と一致し、
学習データが事前学習分布の内側にあるという説明と整合するが、
規模の小ささなど競合する説明を排除できてはいない。

本研究の寄与は、\textbf{「確保されている」と「計算している」を分離して記録する}
という単純な機構にある。この機構がなければ、
16 ランクが 1 ノードに集中した実行も、勾配同期が行われない実行も、
成功として記録されうる。実際、本研究の実装過程では、
コンテナのエントリポイントが何も実行しないまま 42 秒で正常終了し、
ジョブが成功として記録される事象が発生した。ログを開いて初めて発覚したものである。
分散実行の検証は、実行基盤の報告する成否とは独立に行われる必要がある。

\bibliographystyle{plain}
\bibliography{references}

\end{document}
